% =========================================================================
% Capítulo 3: Cinética de sistemas de partículas
% Tema 01: Sistemas de partículas
% Problema 03: Impacto inelástico en péndulo acoplado a carro
% Origen: Ayudantía 29, 2026
% =========================================================================
\subsubsection{Impacto inelástico en péndulo acoplado a carro}

Una partícula $A$ de masa $m/2$ se desplaza con una velocidad inicial $v_0$ hacia la derecha e impacta a una partícula $B$ de masa $m/2$, inicialmente en reposo. La partícula $B$ está conectada a un carro $C$ de masa $m$ mediante una barra rígida sin masa de longitud $l$. El carro se desliza sobre una superficie horizontal sin fricción. Sabiendo que la colisión es perfectamente inelástica, realice lo siguiente:

\begin{enumerate}
    \item Defina los marcos de referencia y sistemas de coordenadas adecuados.
    \item Analice la cinemática del sistema para determinar la velocidad de cada partícula y la velocidad del centro de masa inmediatamente después del impacto.
    \item Determine el máximo ángulo $\theta$ que alcanza la barra después de la colisión utilizando principios de conservación.
\end{enumerate}

\begin{figure}[!ht]
    \centering
    \begin{tikzpicture}[>=Latex, scale=1.4]
        % --- Parámetros ---
        \def\Lbar{2.0}
        \def\Wcar{2.2}
        \def\Hcar{0.6}

        % --- Suelo ---
        \fill[pattern=north east lines] (-3, -0.2) rectangle (3, 0);
        \draw[thick] (-3, 0) -- (3, 0);
        \draw[thin] (-1.0, 0.4) -- (-0.6, 0.1);

        % --- Carro C ---
        \draw[thick, fill=gray!10] (-\Wcar/2, 0.1) rectangle (\Wcar/2, \Hcar+0.1);
        \node at (0.7, \Hcar/2+0.1) {$C$};
        
        \coordinate (O) at (0, \Hcar-0.1);
        \filldraw (O) circle (1.5pt) node[left, xshift=-2pt] {$O$};

        % --- Barra y Partícula B ---
        \draw[line width = 4pt] (O) -- (0, \Hcar-0.1-\Lbar);
        \draw[line width = 3pt, gray!80, line cap = round] (O) -- (0, \Hcar-0.1-\Lbar);
        \shadedraw[ball color=gray!30] (0, \Hcar-0.1-\Lbar) circle (0.18) node[right, xshift=5pt] {$B$};
        \node[left] at (0, \Hcar-0.1-\Lbar/2) {$l$};
        \filldraw (O) circle (0.6pt);
        
        % --- Partícula A ---
        \coordinate (A) at (-2.0, \Hcar-0.1-\Lbar);
        \shadedraw[ball color=gray!30] (A) circle (0.18) node[left, xshift=-5pt] {$A$};
        \draw[->, thick] (A) -- ++(0.7, 0) node[midway, above] {$\mathbf{v}_0$};
        \node[above left, font=\small] at (-0.8, 0.8) {Sin Fricción};

        % --- Gravedad ---
        \draw[->, thick] (2.2, 1.2) -- (2.2, 0.4) node[midway, right] {$\mathbf{g}$};

    \end{tikzpicture}
    \caption{Sistema de impacto inelástico entre una partícula proyectil y un péndulo montado sobre un carro móvil.}
    \label{fig:impacto_carro_pendulo}
\end{figure}


% --- SOLUCIÓN ---------------------------------------------------------------
\begin{solucion}{Impacto inelástico en péndulo acoplado a carro}
\textbf{Planteamiento.} Marco inercial $F$ con base $\{\vec{e}_x,\vec{e}_y,\vec{e}_z\}$ y marco $A$ ligado a la barra $OB$ con base $\{\vec{e}_r,\vec{e}_{\theta},\vec{e}_z\}$. Coordenadas: $x$ (posición horizontal de $O$, medida desde el lugar del choque) y $\theta$ (ángulo de la barra respecto de $\vec{e}_x$). Todas las velocidades se miden en $F$. Se supone que la barra sin masa sólo transmite fuerza en su propia dirección y que, tras el choque, se estudia el sistema $B+C$ (con centro de masa $G$); $A$ interviene sólo mediante el impulso $\hat P$ que le entrega a $B$.
 
\begin{center}
\begin{tikzpicture}[>=Stealth, scale=1.2, font=\small]
  \def\L{2.0}\def\hc{0.9}\def\wc{1.6}\def\yO{0.45}\def\xo{1.4}\def\th{30}
  \pgfmathsetmacro{\xB}{\xo+\L*sin(\th)}
  \pgfmathsetmacro{\yB}{\yO-\L*cos(\th)}
  \pgfmathsetmacro{\yc}{\hc/2}
  \coordinate (O) at (\xo,\yO);
  \coordinate (B) at (\xB,\yB);
  % suelo y carro
  \fill[pattern=north east lines, pattern color=black!55] (-2.2,-0.22) rectangle (4.8,0);
  \draw[thick] (-2.2,0) -- (4.8,0);
  \draw[thick, fill=black!8] (\xo-\wc,0) rectangle (\xo+\wc,\hc);
  \node at (\xo-1.05,\yc) {$m$};
  \node at (\xo+1.15,\yc) {$C$};
  % cota de x
  \draw[dashed, gray] (0,1.5) -- (0,\yO-\L-0.3);
  \draw[dashed, gray] (\xo,\hc) -- (\xo,1.5);
  \draw[<->, gray] (0,1.3) -- (\xo,1.3) node[midway, above] {$x$};
  % barra, vertical de referencia y ángulo
  \draw[dashed, gray] (O) -- ++(0,-\L);
  \draw[->] (B) -- ++(-60:0.9) node[below right] {$\vec{e}_r$};
  \draw[->] (B) -- ++(30:0.9) node[above right] {$\vec{e}_{\theta}$};
  \draw[line width=2.2pt, black!70] (O) -- (B) node[midway, right=2pt] {$l$};
  \fill (O) circle (0.045);
  \node[above left=0pt] at (O) {$O$};
  \draw[<->, gray] ($(O)+(0,-0.9)$) arc[start angle=-90, end angle=-60, radius=0.9];
  \node at ($(O)+(-75:1.25)$) {$\theta$};
  \node[circle, draw, thick, fill=black!8, inner sep=0pt, minimum size=5.5mm] at (B) {$B$};
  % marco F
  \coordinate (F0) at (-1.8,-0.7);
  \draw[->] (F0) -- ++(0.8,0) node[right] {$\vec{e}_y$};
  \draw[->] (F0) -- ++(0,-0.8) node[below] {$\vec{e}_x$};
  \draw (F0) circle (0.07); \fill (F0) circle (0.02);
  \node[above left=-1pt] at (F0) {$\vec{e}_z$};
  \node at ($(F0)+(-0.45,-0.5)$) {$F$};
\end{tikzpicture}
\end{center}
 
\medskip
\textbf{1) Velocidades justo después del choque.}
Choque perfectamente inelástico ($e=0$), con $n$ la dirección normal al plano de colisión ($\vec{e}_y$):
$$e=\frac{v'_{B_n}-v'_{A_n}}{v_{A_n}-v_{B_n}}=0\quad\Rightarrow\quad v'_{B_n}=v'_{A_n}\equiv v'_n .$$
La barra sólo transmite fuerza según $\vec{e}_x$, así que en $\vec{e}_y$ el único impulso sobre $B$ es el de $A$ y se conserva el momentum del par $A$--$B$:
$$m_A\,\vec{v}_{A_n}+\underbrace{m_B\,\vec{v}_{B_n}}_{0}=m_A\,\vec{v}\,'_{A_n}+m_B\,\vec{v}\,'_{B_n}$$
$$\frac m2\,(v_0\vec{e}_y)=\frac m2\,\vec{v}\,'_n+\frac m2\,\vec{v}\,'_n
\;\Rightarrow\; \vec{v}\,'_A=\vec{v}\,'_B=\frac{v_0}{2}\,\vec{e}_y$$
 
\textbf{Impulso sobre $B$.} Como $B$ parte del reposo:
$$\hat P=m_B\,\vec{v}\,'_B=\frac m2\cdot\frac{v_0}{2}\,\vec{e}_y=\frac{m v_0}{4}\,\vec{e}_y .$$
Este impulso produce un movimiento lineal del centro de masa y otro rotacional de la barra.
 
\textbf{Centro de masa de $B+C$.} El centro de masa estaba en reposo antes del choque:
$$\underbrace{(m_B+m_C)\,\vec{v}_G}_{0}+\hat P=(m_B+m_C)\,\vec{v}_G'
\;\Rightarrow\;
\vec{v}_G'=\frac{m v_0/4}{m/2+m}=\frac{m v_0/4}{\tfrac32 m}=\frac{v_0}{6}\,\vec{e}_y .$$
 
\textbf{Velocidad del carro y giro de la barra.} Por cinemática, con $\vec{e}_{\theta}=\vec{e}_y$ en $\theta=0$:
$$\vec{v}_C=\dot x\,\vec{e}_y,\qquad \vec{v}_B=\dot x\,\vec{e}_y+l\dot\theta\,\vec{e}_{\theta} .$$
Con el subíndice $i$ para el instante justo después del choque, hay dos condiciones. La del centro de masa:
$$ \vec{v}_G'=\frac{m_C\,\dot x_i+m_B\,(\dot x_i+l\dot\theta_i)}{m_C+m_B}\,\vec{e}_y=\Big[\dot x_i+\frac13\,l\dot\theta_i\Big]\vec{e}_y=\frac{v_0}{6}\,\vec{e}_y
\;\Rightarrow\; 6\,\dot x_i+2\,l\dot\theta_i=v_0 ,$$
y la velocidad de $B$ obtenida en el choque, $\dot x_i+l\dot\theta_i=v_0/2$, es decir, $2\,\dot x_i+2\,l\dot\theta_i=v_0$. Resolviendo:
$$\dot x_i=0,\qquad \dot\theta_i=\frac{v_0}{2\,l}.$$
\respuesta{$\displaystyle \vec{v}\,'_A=\vec{v}\,'_B=\frac{v_0}{2}\,\vec{e}_y,\qquad \vec{v}\,'_C=\vec{0},\qquad \vec{v}_G'=\frac{v_0}{6}\,\vec{e}_y$}
 
\medskip
\textbf{2) Máximo ángulo de la barra.}
Tras el choque no actúan fuerzas no conservativas (sin roce, la normal no trabaja y la barra rígida sin masa no hace trabajo neto), luego se conserva la energía de $B+C$: ${}^{F}T_f+{}^{F}U_f={}^{F}T_i+{}^{F}U_i$, con estado inicial justo después del choque y final en $\theta_{\max}$.
 
\textbf{Energía potencial} (referencia $U=0$ a la altura de $O$; la del carro es constante):
$${}^{F}U_i=-\frac m2\,g\,l,\qquad {}^{F}U_f=-\frac m2\,g\,l\cos\theta_{\max}.$$
 
\textbf{Energía cinética (König).} ${}^{F}T={}^{F}T^{(1)}+{}^{F}T^{(2)}$, con
$${}^{F}T^{(1)}=\tfrac12\,(m_B+m_C)\,\vec{v}_G\cdot\vec{v}_G,\qquad
{}^{F}T^{(2)}=\tfrac12\sum_k m_k\,(\vec{v}_k-\vec{v}_G)\cdot(\vec{v}_k-\vec{v}_G).$$
\begin{itemize}
  \item \emph{Final:} en $\theta_{\max}$, $\dot\theta=0$ y las partículas no se mueven respecto del centro de masa, así que ${}^{F}T_f^{(2)}=0$.
  \item \emph{Centro de masa:} no hay fuerzas horizontales sobre $B+C$, luego $\vec{v}_{G,f}\cdot\vec{e}_y=\vec{v}_{G,i}\cdot\vec{e}_y=v_0/6$ y ${}^{F}T_i^{(1)}={}^{F}T_f^{(1)}$.
  \item \emph{Inicial:} con $\vec{v}_{C,i}=\vec{0}$, $\vec{v}_{B,i}=\tfrac{v_0}{2}\vec{e}_y$, $\vec{v}_{G,i}=\tfrac{v_0}{6}\vec{e}_y$:
\end{itemize}
$${}^{F}T_i^{(2)}=\frac12\Big[m\Big(0-\frac{v_0}{6}\Big)^2+\frac m2\Big(\frac{v_0}{2}-\frac{v_0}{6}\Big)^2\Big]
=\frac12\Big[\frac{m v_0^2}{36}+\frac{m v_0^2}{18}\Big]=\frac{m v_0^2}{24}.$$
 
\textbf{Balance.} Como ${}^{F}T^{(1)}$ no cambia y ${}^{F}T_f^{(2)}=0$, queda $U_f={}^{F}T_i^{(2)}+U_i$:
$$-\frac m2\,g\,l\cos\theta_{\max}=\frac{m v_0^2}{24}-\frac m2\,g\,l
\;\Rightarrow\;
\cos\theta_{\max}=1-\frac{v_0^2}{12\,g\,l}$$
\respuesta{$\displaystyle \theta_{\max}=\arccos\!\left(1-\frac{v_0^2}{12\,g\,l}\right)$}
Válido mientras el argumento sea mayor que $-1$, es decir, $v_0^2<24\,g\,l$.
 
\medskip
\textbf{Verificación.} (i) Momentum total en $\vec{e}_y$: antes, $m_A v_0=\tfrac12 m v_0$; después, $\tfrac m2\tfrac{v_0}{2}+\tfrac m2\tfrac{v_0}{2}+0=\tfrac12 m v_0$. (ii) Sin König: en $\theta_{\max}$, $B$ y $C$ van a la misma velocidad $u$ y $\tfrac32 m\,u=\tfrac m2\tfrac{v_0}{2}$, así $u=v_0/6$; la energía da $\tfrac{m v_0^2}{16}=\tfrac{m v_0^2}{48}+\tfrac m2 g l(1-\cos\theta_{\max})$, que lleva al mismo resultado. (iii) Si $v_0\to0$, $\theta_{\max}\to0$.
 

\patoojo{Durante el choque, que dura un instante, la barra solo puede transmitir fuerza a lo largo de su propia dirección. Como el impulso de $A$ sobre $B$ es perpendicular a la barra, el carro $C$ no recibe impulso horizontal y por eso sale del choque con $\dot{x}=0$. Toda la velocidad de $B$ se la lleva la rotación de la barra.}

\patotruco{Para el ángulo máximo no hace falta resolver ninguna ecuación diferencial: en $\theta_{\max}$ la barra deja de girar y $B$ y $C$ se mueven juntos con la misma velocidad $u$. Conservación de momentum horizontal ($u=v_0/6$) más conservación de energía, y listo.}
\end{solucion}
